\problem{P10}{Dependent Binary Choices}{Markov correlations}
Starting from the unit interval, choose its left or right half at step $j$ according as $B_j=0$ or $1$. The first bit is fair. Thereafter the bit stays unchanged with probability $p\in[0,1]$ and flips with probability $1-p$, according to a Markov transition. Let $X$ be the limiting position. Find its mean and variance, describe its law for $p=0,1/2,1$, and bound the error from using the first $n$ bits.
\solution The nested interval lengths tend to zero, and their unique limiting point is
\[
X=\sum_{j\ge1}2^{-j}B_j.
\]
The chain is stationary with $\E B_j=1/2$. Put $Z_j=B_j-1/2$ and $\lambda=2p-1$. Since $\E(Z_{j+1}\mid Z_j)=\lambda Z_j$, iteration gives
\[
\Cov(B_i,B_j)=\frac14\lambda^{|i-j|}.
\]
Absolute convergence permits rearrangement. Using $\sum_i4^{-i}=1/3$,
\[
\begin{aligned}
\Var X
&=\frac14\sum_i4^{-i}
+\frac12\sum_{k\ge1}\lambda^k\sum_{i\ge1}2^{-i}2^{-(i+k)}\\
&=\frac1{12}+\frac16\sum_{k\ge1}(\lambda/2)^k
=\boxed{\frac{2+\lambda}{12(2-\lambda)}}.
\end{aligned}
\]
Also $\boxed{\E X=1/2}$.

For $p=0$, the bits alternate, giving $X=1/3,2/3$ with equal probabilities. For $p=1$, all bits agree, giving $X=0,1$ equally. For $p=1/2$, the bits are independent; every binary interval of length $2^{-n}$ has probability $2^{-n}$, proving that $X$ is uniform. The corresponding variances are $1/36,1/12,1/4$.

For $X_n=\sum_{j=1}^n2^{-j}B_j$, we always have $0\le X-X_n\le2^{-n}$. Using the midpoint of the selected interval reduces the absolute error bound to $2^{-n-1}$.
\insight Fair marginal bits do not ensure a uniform limiting point. Dependence between digits accumulates through the geometric weights and changes the variance.
